\section{Capture, edge charges, and snapshots}\label{sec:capture}

The frequency lists in this section serve two different purposes. A cumulative
list gives a deterministic approximation by signed sums. At selected spatial
prefixes we compress the relevant signed sums into a smaller list of
representatives. The first list controls how long approximation can fail; only
the smaller lists will enter the probability estimates of
\Cref{sec:paths}. All constructions take place in the finite procedure of
\Cref{sec:procedure}, uniformly in its depth and in any frequency cutoff.

\subsection{The two laws on a fixed plateau}

For static pieces, start a plateau at entry to static handling and restart
whenever the incoming mean exceeds twice its current normalization. There are
no membership tests on these plateaus. They obey the same common stopping rule
as live pieces. Discard zero pieces.

Fix a live or static plateau $e$, with root $p_e$ and normalization
$\alpha=\alpha_e>0$. Under $\mathbf U_e$, spatial uniform probability on
$p_e$, follow only this plateau's surviving density, through its failures,
until its exit or the common stop. Its jumps on an output path occur at
consecutive depths. At a jumping node $p$, write $a$ for the final
\emph{post} density, after all removals there, and $m_p=\mathbb E_p a$.
The two prospective child measures use this same density restricted to the
children, \emph{before} any processing in a child. Thus
\begin{equation}\label{eq:07-mass-bounds}
  m_p/\alpha\le2,
  \qquad \mathbb E_{p'}a/\alpha\le4
  \quad(p'\text{ a child of }p).
\end{equation}
Let $\xi_p$ be the mass removed at $p$, measured as its local spatial average
divided by $\alpha$. Full removal at an exit can be included. The removed
input sets are disjoint, so
\begin{equation}\label{eq:07-kill-bounds}
 \mathbb E_{\mathbf U_e}\sum_p\xi_p\le1,
 \qquad
 \mathbb E_{\mathbf U_{p_0}}\sum_{p\text{ after the post state }p_0}\xi_p
       \le m_{p_0}/\alpha\le2.
\end{equation}
The second sum excludes removals already made at $p_0$. Here and below a sum
along a path includes only states of the specified plateau. Put
\[
 d_1=d\quad\text{on a live plateau},\qquad
 d_1=d+2+\log^-\alpha\quad\text{on a static plateau}.
\]

For comparison, let $\mathbf P_e$ be the input experiment starting with the
entry piece normalized to probability on $p_e$. It observes the spatial bits
of its input $Y$ and its actual memberships; hence it can leave the plateau
on a success. It may subsequently reveal more spatial bits, with no further
observations, when a finite extension is useful. A final jumping post state
has probability
\begin{equation}\label{eq:07-state-conversion}
 \mathbf P_e(\text{post state at }p)
       =\mathbf U_e(p)\frac{m_p}{\alpha}.
\end{equation}
Conditional on that state, the input density is $a/m_p$ relative to uniform
on $p$. If the experiment instead begins at a jumping post state $p_0$, the
corresponding probability of a later post state $p\subseteq p_0$ is
$\mathbf U_{p_0}(p)m_p/m_{p_0}$. These identities follow directly because
each survivor is the entry density restricted by its spatial prefix and
membership events.

The uniform output law does not condition on a random input membership. At
each reached output prefix, the survivor and all removals are already
determined functions. Both prospective children can therefore be examined
before the next fair output bit. For a live plateau the side information
will be its fixed main label $Z$, with conditional density $K_A(z-Y)$ given
$Y$. The discrete decisions do not inspect its realized value. For a scalar
construction there is no side information.

\begin{lemma}[Local entropy and Fisher budgets]\label{lemma:07-local-budget}
On the plateau just described, the expected sum of spatial bit relative
entropies to a fair bit is $O(d_1)$, with or without the fixed side label.
This includes every finite extension after exit. If finite-valued
observations are made at predictable states and their total conditional
entropy cost has expectation at most $B$, the bound becomes $C d_1+B$.

Let $M$ be the normalized joint post mass measure of the side information and
the finite observations so far, and let $D_{\rm jump}$ be the half-difference
of its prospective child measures. Then
\begin{equation}\label{eq:07-local-fisher}
 \mathbb E_{\mathbf U_e}\sum_p
       \int\frac{|D_{\rm jump}|^2}{M}\le C(d_1+B).
\end{equation}
For a forward estimate from a live jumping post state $p_0$, restart the
input experiment with its post density, retaining the fixed main label
when present but discarding all earlier auxiliary observations. The same
bounds then hold with $d_1=d$, with $B$ charging the auxiliary observations
made in this restarted experiment, and with an extra factor
$m_{p_0}/\alpha\le2$ in the mass conversion. Already completed removals
are excluded. In particular the future main-label Fisher sum has
conditional expectation at most $Cd$. If earlier auxiliary information
$W$ is retained instead, its conditional information
$\mathrm I(Y;W\mid Z)$ under the restarted input law must also be added
to the initial budget; omit $Z$ when there is no side label.
All quotients are zero at zero mass; for measures, densities to a common
dominating measure are understood.
\end{lemma}

\begin{proof}
The initial density deficit relative to spatial uniform is at most
$d-\log\alpha\le C d_1$. On a live plateau with jumps,
$\alpha\ge e^{-d}$. To pay for its fixed label, let $p_*$ be the interval
of the last main addition and $c_*$ its center. Compare its channel with the
fixed density $K_{S_{p_*}}(z-c_*)$. Writing $Z=Y+\theta$ gives
\[
 \begin{split}
 \mathrm I(Y;Z)&\le\mathbb E_Y
    D\bigl(K_A(z-Y)\,dz\,\|\,K_{S_{p_*}}(z-c_*)\,dz\bigr)\\
 &=F(A)-F(S_{p_*})-\mathbb E A(\theta)
       +\mathbb E S_{p_*}(Y+\theta-c_*)\\
 &\le F(A)-\log R_{p_*}+C\le Cd.
 \end{split}
\]
Indeed $p_e\subseteq p_*$, the channel includes $S_{p_*}$, and the spatial
variance and form monotonicity in \Cref{sec:thermal} bound the last energy
by a constant: use the phase triangle inequality and
$|Y-c_*|\le |p_*|/2$. The negative energy can be dropped. The initial
dummy channel has zero information. This estimate holds for every input
law supported on $p_e$, so also for the post law at a later live jumping
state.

Within one live plateau a true path has at most one success, which is
terminal for that plateau. Its success probability at a test is at least
$c b_0^2$, so its success surprisal is $O(1+\log h)$. The accumulated
failure surprisal has an exponential tail by the failure-stretch
supermartingale used in \Cref{sec:procedure}. Thus the expected total
membership entropy cost is $O(1+\log h)$. Static plateaus have no such cost.
Forward from a live post state, $m_p\ge e^{-d}$, so the initial density
deficit is again at most $2d$, and the same membership argument applies.
This forward argument begins without earlier auxiliary observations.
Retaining such information $W$ instead increases the initial information
term from $\mathrm I(Y;Z)$ to
$\mathrm I(Y;Z,W)=\mathrm I(Y;Z)+\mathrm I(Y;W\mid Z)$.

For completeness, apply the entropy chain rule to any finite remaining
spatial bit string. Its expected deficit relative to the uniform string
starts below the density deficit, plus the information of any label already
observed. Revealing a bit consumes its conditional deficit relative to a
fair bit. An additional observation increases the expected remaining
deficit by its conditional information about the remaining string, at most
its conditional Shannon entropy if finite-valued. The remaining deficit is
nonnegative at stopping. Membership entropy can be bounded before
conditioning on side and auxiliary observations, using the discrete state
history that prescribes the tests. This proves the first assertion,
including finite extensions.

At a post state $p$, the joint observations have conditional density
$(\alpha/m_p)M$ under truth and the next bit bias is $D_{\rm jump}/M$.
Their averaged squared bias is therefore
\[
 \frac{\alpha}{m_p}\int\frac{|D_{\rm jump}|^2}{M}.
\]
Observing further past information only increases the mean squared
conditional bias, by conditional Jensen. For a bit with bias $\theta$,
its relative entropy to fair is at least $\theta^2/2$. Multiply by
\Cref{eq:07-state-conversion} and sum: the factors $m_p/\alpha$ and
$\alpha/m_p$ cancel. This proves \Cref{eq:07-local-fisher} without any
lower bound on surviving mass. Starting at $p_0$ instead leaves the factor
$m_{p_0}/\alpha\le2$, as asserted.
\end{proof}

\subsection{The capture statement}

There are two separate constructions at each allowed dyadic level $b$.
For the \emph{main} criterion, used on live plateaus for $b_0\le b\le1$,
set $\kappa=2$ and say that $u$ is high if
$\max_{p'}G_{p'}^u\ge b^2$, using the prospective labeled child moments.
For the \emph{scalar} criterion, set $\kappa=1$ and say that $u$ is high if
$\max_{p'}|f_{p'}^u|\ge b$, where
$f_{p'}^u=\mathbb E_{p'}[a\chi_u]/\alpha$ has no label. Scalar levels are
$0<b\le b_0$ on live plateaus and $0<b\le1$ on static plateaus. In
particular the main and scalar constructions at $b_0$ are distinct.

Set
\begin{equation}\label{eq:07-edge-definition}
 l_b=\left\lceil20\log_2\frac{d_1+2}{b}\right\rceil,
 \quad \delta_p=2^{-l_b}|p|,
 \quad e_{p,b}=\frac{\mathbb E_p[a\mathbf1_{\mathrm{edge}(p,b)}]}{\alpha},
\end{equation}
where $\mathrm{edge}(p,b)$ consists of points inside $p$ within
$4\delta_p$ of an endpoint or the midpoint. The jump is \emph{bad} for
this criterion and level if $e_{p,b}\ge c b^\kappa$, for a sufficiently
small fixed positive $c$. Everything in this definition is known at the
post state before the next output bit.

\begin{proposition}[Capture on a plateau]\label{prop:capture}
In the finite procedure, for every plateau, allowed level, and criterion
above, there is a cumulative finite list $D$ of integers, initially empty,
whose appends are prescribed at jumping post states, such that, after the
appends at $p$,
\begin{equation}\label{eq:source-14}
 \begin{gathered}
 u\text{ high and the jump not bad}\quad\Longrightarrow\quad
 |u-v||p|\le\poly(1+|D|,d_1,1/b)
       \quad\text{for some }v\in[D],\\
 [D]=\left\{\sum_{n\in D}\eta_n n:\eta_n\in\{0,1,-1\}\right\},
 \qquad \mathbb E_{\mathbf U_e}|D|_{\rm final}\le C d_1 b^{-C}.
 \end{gathered}
\end{equation}
The construction is independent of the caps used below.

The bad-edge charge at $p$ and level $b$ is
$e_{p,b}\mathbf1_{\{p\text{ bad at level }b\}}$. Its global live sum,
with weights $\mathbf U(p)\alpha_e$, is $O(d)$ when one scalar level is
selected predictably at each piece, and
$O(\poly(h)(h+\log d))$ when all main levels are summed. The local
one-scalar-level charge is $O(d_1)$ under $\mathbf U_e$. Forward from a
live post state the conditional expected number of bad jumps at a fixed
main level is at most $d\poly(h)$.

For each deterministic dyadic $D_{\max}\ge1$, the scheme stopped when
$|D|>D_{\max}$ admits snapshots with expected number
$\poly(D_{\max},d_1,1/b)$. A snapshot at level $b$ stores at most
$\lceil e^{k_{b/2}}\rceil$ main representatives, or
$\exp(C(d_1+\log(1/b)))$ scalar representatives. Each birth and its list
are determined by the spatial prefix before its birth jump.

At a jump $p$, call a high integer $u$ activated when some $v\in[D]$
satisfies $|u-v||p|\le\sigma$ in the main case, or
$|u-v||p|\le c_8b^4$ in the scalar case, for a sufficiently small fixed
$c_8>0$. At activation its chosen representative $s$
has, for the remainder of that plateau, a fixed decomposition
\[
 \begin{array}{ll}
 u=s+n+a',\quad \Delta_A(n)\le r_{\rm sep},\quad
          |a'||p_{\rm act}|\le C\sigma,&\text{main},\\
 u=s+a',\quad |a'||p_{\rm act}|\le c_1b^4,&\text{scalar},
 \end{array}
\]
where $p_{\rm act}$ is the activation node and $c_1$ can be made a
sufficiently small absolute constant. Before first activation, all
non-bad high jumps for this $u$ lie in a single block of at most
\[
 H_b=C\log(CD_{\max}d_1/b)+C\log(1/\sigma)
\]
consecutive depths, as long as the cap is not breached. The construction
does not use sampled side or auxiliary values on an output path.
\end{proposition}

We prove the assertions in turn. The intermediate lemmas also specify the
conditioning needed for their later use.

\subsection{Broad capture by trimmed masks and polynomial products}

Suppose $u$ is high at a jumping post piece. Polarize on one high child:
there is a measurable complex $v(z)$ with $|v(z)|\le1$ (a constant in the
scalar case) such that the normalized real bias
\begin{equation}\label{eq:07-polarized-bias}
 \operatorname{Re}\frac1\alpha\mathbb E_p
       [a\mathbf1_{p'}(Y)v(Z)\chi_u(Y)]\ge b^\kappa/2.
\end{equation}
For the main criterion, $|f_{p'}^u|\le M_{p'}$ implies
$\int|f_{p'}^u|\ge G_{p'}^u$; choosing its pointwise conjugate phase proves
the assertion. In particular every append state below satisfies
$m_p/\alpha\ge b^\kappa/2$.

We first replace the child indicator by a trigonometric polynomial. Write
$n=|D|$ at this stage. Trim a distance $\delta_p$ from both ends of $p'$,
and convolve its indicator on the unit circle with the positive normalized
kernel proportional to
\[
 \left(\frac{\sin(\pi N x)}{N\sin(\pi x)}\right)^{2s},
 \qquad
 N=\left\lceil C2^{l_b}/|p|\right\rceil,
 \qquad
 s=\left\lceil C(1+n+\log(1/b))\right\rceil.
\]
Its continuous values at integers are understood. The result $P_p$ is a
real trigonometric polynomial, $0\le P_p\le1$, of degree at most $sN$.
The unnormalized kernel integral is at least $c/(N\sqrt{s})$: on circle
distance at most $c/(N\sqrt{s})$ from zero, the ratio has $2s$-th power
bounded below by a positive numerical constant. The bound
$|\sin\pi x|\ge2\operatorname{dist}(x,\mathbb Z)$ gives its unnormalized
tail outside distance $\delta_p$ at most $Ce^{-c's}/N$, after choosing the
constant in $N$ large. Consequently the normalized tail is at most
$C\sqrt{s}e^{-c's}$. Outside the edge inside $p$, the mask approximates
$\mathbf1_{p'}$ uniformly with this error. Integrating the convolution in
the other order also gives
\begin{equation}\label{eq:07-mask-bounds}
 \deg P_p\le \frac{C2^{l_b}(1+n+\log(1/b))}{|p|},
 \qquad
 \int_{\mathbb T\setminus p}P_p
 \le |p|\exp[-C'(1+n+\log(1/b))].
\end{equation}
Here $C'$ can be chosen as large as needed by increasing the fixed kernel
constants. The edge error in \Cref{eq:07-polarized-bias} is at most
$e_{p,b}$, and the remaining error is at most
$2C\sqrt{s}e^{-c's}$. Thus, at a non-bad jump, the replacement retains bias
$\ge c_2b^\kappa$. This argument also covers intervals that cross the
circle coordinate boundary.

For each fixed side value $z$, keep a full-circle polynomial product $Q_z$,
initially one. Its restriction to the current $p$, normalized by its
integral there, is the reference probability density for $Y$ conditional on
that side value. After $n$ appends its factors lie between $1/2$ and $2$,
so
\[
 2^{-n}\le Q_z\le2^n,\qquad \int_pQ_z\ge2^{-n}|p|.
\]
If a non-bad high $u$ is outside $[D]$ enlarged by the sum of all previous
mask bandwidths and the prospective bandwidth in
\Cref{eq:07-mask-bounds}, append it and multiply $Q_z$ by
\begin{equation}\label{eq:07-product-factor}
 1+\eta b^\kappa\operatorname{Re}[v(z)\chi_u P_p],
\end{equation}
where $\eta>0$ is a sufficiently small fixed constant. The old product has
Fourier support inside $[D]$ enlarged by its accumulated mask bandwidths.
Thus the full-circle integral of $Q_z\chi_u P_p$ vanishes. Restriction to
$p$ changes its normalized integral by at most
\[
 \frac{2^n\int_{\mathbb T\setminus p}P_p}{2^{-n}|p|}
 \le4^n\exp[-C'(1+n+\log(1/b))].
\]
Choose $C'$ so this is a suitably small multiple of $b^\kappa$, uniformly
in $n$ and $z$.

The relevant entropy is the true average, over the current conditional side
law, of the relative entropy of $Y$ conditional on that side value to this
reference. Its drop at an append equals the true expectation of the
logarithm of \Cref{eq:07-product-factor}, minus the average logarithm of
the reference expectation of that factor. By
\Cref{eq:07-mass-bounds}, the retained normalized bias becomes at least
$c_2b^\kappa/2$ when divided by $m_p/\alpha$ to obtain a probability-law
bias. The inequality $\log(1+x)\ge x-Cx^2$ for $|x|\le1/2$ therefore
bounds the first logarithm below by $c\eta b^{2\kappa}$, after fixing
$\eta$ small. The leakage estimate bounds the normalization logarithm by
a smaller multiple of the same quantity. Every append hence consumes at
least $c_3b^{2\kappa}$ of this conditional entropy.

Here is the budget across states, including the change in the side law.
Denote a discrete state history by $H$, and the current reference by
$Q_{H,z}$, restricted and normalized on its interval. Set
\[
 \mathcal B(H)=\mathbb E_{Z\mid H}
       D(\mathcal L(Y\mid H,Z)\|Q_{H,Z}).
\]
At a spatial bit $X$, restrict the reference to the selected child. The
relative entropy chain rule gives exactly
\[
 \sum_x\mathbb P(x\mid H)\mathcal B(H,x)
 =\mathcal B(H)-\mathbb E_{Z\mid H}
       D(\mathcal L(X\mid H,Z)\|\mathcal L_{Q_{H,Z}}(X))
 \le\mathcal B(H).
\]
At a membership observation $C$, keep the old reference on both outcomes.
Its average entropy increment is
$\mathrm I(Y;C\mid H,Z)\le\mathrm H(C\mid H)$, since the current side posterior is used
on each outcome. Nonnegative entropy on exited branches can be discarded.
Initially the budget is
$D(\mathcal L(Y)\|\mathbf U_e)+\mathrm I(Y;Z)\le Cd_1$; the membership costs are
those in \Cref{lemma:07-local-budget}. Thus
\begin{equation}\label{eq:07-append-budget}
 \mathbb E_{\mathbf P_e}(\text{number of appends})
       \le Cd_1b^{-2\kappa},\qquad
 \mathbb E_{\mathbf U_e}|D|_{\rm final}
       \le Cd_1b^{-3\kappa}.
\end{equation}
The second estimate uses \Cref{eq:07-state-conversion} and
$m_p/\alpha\ge b^\kappa/2$ at every append. At a fixed positive-mass
state the entropy is finite and drops strictly, so the append loop is
finite. This finiteness propagates down every positive-probability branch
of the finite processing tree.

If no further append is possible, every non-bad high $u$ is within the
accumulated bandwidth of $[D]$. Previous append intervals are ancestors of
$p$, and previous list lengths are at most the current length. Hence that
bandwidth, multiplied by $|p|$, is bounded by
\[
 C2^{l_b}(n+1)(n+1+\log(1/b)).
\]
For example $C(d_1+2)^{20}b^{-21}(n+1)^2$ is a permissible fixed
polynomial upper bound. This proves \Cref{eq:source-14}. The signs,
products, and appended frequencies are chosen from the laws at each state;
no realized side value is used to decide an append.

\Needspace{11\baselineskip}
\subsection{Charging thin spatial edges}

\begin{lemma}[Edge charges]\label{lemma:07-edge-charges}
With the bad-jump convention \Cref{eq:07-edge-definition}, the following
bounds hold.
\begin{enumerate}[label=\textup{(\roman*)}]
\item Across all live jumping pieces, the sum of
$\mathbf U(p)\alpha_e e_{p,b}\mathbf1_{\rm bad}$ is at most $Cd$ if a
single scalar level $b$ is selected at each piece before its next output
bit, without using noisy label values.
\item The corresponding sum over all main levels and live pieces is at
most $C\poly(h)(h+\log d)$.
\item Within one plateau, the expected one-scalar-level sum under
$\mathbf U_e$ is at most $Cd_1$.
\item From any live jumping post state, under spatial uniform conditional
on that prefix, the expected number of future bad jumps of the same plateau
at main level $b$ is at most $Cd b^{-2}$, and hence $d\poly(h)$.
\end{enumerate}
\end{lemma}

\begin{proof}
For (i) and (ii), fix the global unlabelled input experiment of
\Cref{prop:live-budgets}, following memberships through all live epochs
and freezing the terminal piece only at live termination. For (iii) and
(iv), first fix one plateau and use the corresponding local input experiment
of \Cref{lemma:07-local-budget}, freezing the terminal piece at that
plateau's exit. In each application, all windows below belong to one common
finite spatial extension of this fixed experiment and use the same
predictable spatial entropy sequence.

At a final post state the conditional input edge probability is
\[
 q=\frac{\alpha e_{p,b}}{m_p}\ge c b^\kappa/2
 \quad\text{if the jump is bad}.
\]
From this state compare the next $l=l_b$ spatial bits, starting with its
jump bit, to a reference experiment with fair spatial bits. At intervening
membership tests use the very same state-conditional outcome probabilities
as in the true unlabelled experiment. Thus membership likelihood factors
cancel. This reference can be generated by sampling each bit fairly,
regardless of earlier memberships, and by tossing the stated membership
coins on their recorded branches. In particular its spatial path remains
uniform. Continue bits after exits or stops as needed, keeping the piece
fixed and taking no new observations. Histories of zero true mass can be
extended arbitrarily under the reference while leaving its bits fair.

The edge is determined by the next $l$ bits, up to null endpoints, and its
reference probability is at most $16\,2^{-l}$. If $L$ is the
true-to-reference likelihood ratio of this segment, then
\[
 \mathbb P_{\rm true}(\mathrm{edge}\cap\{L<2^{l/2}\}\mid\text{post state})
 \le16\,2^{-l/2}\le q/2.
\]
The last inequality follows uniformly from the choice of $l_b$, the badness
lower bound on $q$, and sufficiently large $d$. Call the starting depth
\emph{marked} if the edge occurs and $\log L\ge(l/2)\log2$.
Conditional on a bad post state, its probability is at least $q/2$.

Let $s_j$ be the predictable relative entropy to fair of the spatial bit at
depth $j$, in this unlabelled true history. If its bias is $\theta_j$,
then $s_j\ge\theta_j^2/2$, whereas the positive part of its realized log
likelihood increment satisfies
\[
 \bigl[\log(1\pm\theta_j)\bigr]_+
       \le |\theta_j|\le\sqrt{2s_j}.
\]
Memberships contribute no log increment to $L$. On a mark, Cauchy--Schwarz
on the length-$l$ window consequently yields
\begin{equation}\label{eq:07-marked-window}
 \sum_{j\text{ in its window}}s_j\ge c_4 l.
\end{equation}
There is at most one final jumping piece at each spatial depth along an
actual input path, although membership splitting may have preceded it.
For any finite family of marked starting indices, select disjoint windows
whose triples cover the union, for example by successively taking a longest
remaining window. The number of starting indices is no larger than a
constant times the sum of the selected window lengths. By
\Cref{eq:07-marked-window}, this is at most $C\sum_j s_j$, since the
selected windows are disjoint. This proves the interval charging bound for
variable window lengths as well as for fixed $l$.

Globally, the unlabelled full spatial entropy budget of
\Cref{prop:live-budgets} is $O(d)$ when live pieces occur. Apply the
interval bound with the one selected scalar level per index. The marking
probability and the state law imply
\[
 \sum_{\rm live\ pieces}\mathbf U(p)\alpha_e
       e_{p,b}\mathbf1_{\rm bad}
 =m_{\rm top}\mathbb E_{\mathbf P}
       \sum_{\rm live\ post\ states}q\mathbf1_{\rm bad}
 \le2m_{\rm top}\mathbb E_{\mathbf P}(\text{number of marks})
 \le Cd.
\]
Only a finite spatial extension is needed for any finite selection of
windows; the bound for the full sum follows by nonnegativity. This also
handles arbitrarily small selected scalar levels.

For a fixed main level, $l_b$ is constant along the live path. Windows
wholly before the common cutoff use only its $O(h)$ truncated unlabelled
bit budget. At most $l_b$ starting indices straddle that cutoff. The
truncated budget remains valid on paths that transferred earlier: keep
their membership state fixed and continue their spatial sampling. Therefore
the main-level weighted edge sum is $O(h+l_b)$. There are $O(\log h)$ main
levels and $l_b=O(\log d+\log(1/b_0))$ on them, proving (ii).

For (iii), use $\mathbf P_e$ and \Cref{lemma:07-local-budget}, with no side
label, in the same proof. The state conversion now has root mass
normalization one, giving $Cd_1$. Starting instead at a live post state
$p_0$, the conversion leaves a factor $m_{p_0}/\alpha\le2$, and the
forward entropy budget is $Cd$. Thus the expected forward sum of
$e_{p,b}\mathbf1_{\rm bad}$ under $\mathbf U_{p_0}$ is at most $Cd$.
At a bad main jump $e_{p,b}\ge cb^2$; dividing gives (iv).
\end{proof}

\subsection{Snapshots in fixed rounded coordinates}

Fix a deterministic dyadic cap $D_{\max}\ge1$ and stop this level's snapshot
scheme on a branch as soon as its independently constructed broad list has
$|D|>D_{\max}$. Use a fixed partition of the unit circle into
$O(D_{\max}/b^6)$ arcs, rounding each $\chi_n(Y)$, $n\in D$, to a chosen
point of its arc with error at most $c_5b^6/D_{\max}$. Track the joint
\emph{finite mass measure} of these rounded coordinates and the side
information, normalized by $\alpha$; do not renormalize it to probability.
For a finite signed measure $\mu$ on $\Omega$, write
$\|\mu\|_{\rm var}=|\mu|(\Omega)$ for its full variation norm,
without a factor $1/2$.

After the broad-list step at each jumping post node, refresh the snapshot
once if it is the first snapshot, if any coordinate was appended, or if
some prospective child pre measure differs in total variation by more than
$c_6b^6$ from the stored \emph{post} measure of the current snapshot. The
last comparison uses the same coordinates on both measures; an append
itself causes replacement. On refresh store the present post measure.
All these decisions inspect deterministic measures at the current prefix.

At a snapshot take a maximal separated subcollection of the elements of
$[D]$ high at threshold $b/2$ at this node. In the main case separation means
that differences avoid
\[
 B_A(r_{\rm sep})+\{j\in\mathbb Z:|j||p|\le\sigma\};
\]
in the scalar case it means that differences have magnitude greater than
$c_7b^4/|p|$. Take all the constants $c_5,c_6,c_7$ sufficiently small.

\begin{lemma}[Snapshot covers and birth count]\label{lemma:07-snapshots}
The main and scalar snapshot lists have respectively the cardinality bounds
in \Cref{prop:capture}. While a snapshot is current, it covers every
$v\in[D]$ that is high at threshold $0.9b$ at a jumping node, with its
stored separation neighborhoods measured at the snapshot's birth interval.
Moreover, if $N_{\rm snap}$ is the number of snapshots on the capped
uniform output path, then
\begin{equation}\label{eq:07-snapshot-count}
 \mathbb E_{\mathbf U_e}N_{\rm snap}
 \le 1+D_{\max}+C b^{-6}
       +C b^{-12}\bigl(d_1+D_{\max}\log(CD_{\max}/b)\bigr).
\end{equation}
Every snapshot is born before its node's next fair bit. Its stored
representatives may subsequently be used on the same plateau even after
replacement of that snapshot.
\end{lemma}

\begin{proof}[Proof of the cover assertions]
At a stabilized live node, a separated high collection of size
$\lceil e^{k_{b/2}}\rceil$ would request another update. The computation
levels include $b/2$. Thus a maximal collection has at most this size and
covers all high elements of $[D]$ by the indicated symmetric neighborhoods.

For scalar lists, Parseval applied to $a\mathbf1_{p'}$ on each prospective
child gives
\[
 \sum_{u\in\mathbb Z}|f_{p'}^u|^2
 =\frac{1}{\alpha^2|p'|^2}\int_{p'}a^2
 \le\frac{C e^d}{\alpha|p|}.
\]
After rotating by the phase at a point of $p'$, its moments vary by at most
$C|u-u'||p|$. Around each separated representative take an integer
neighborhood of radius $c_7b^4/(3|p|)$. These neighborhoods are disjoint;
each contains at least $c b^4/|p|$ integers, with the bound also true when
that number is less than one. On one of the two children their moments are
still at least $c b$. Summing their squared moments over the two children
bounds the list length by
$C e^d/(\alpha b^6)\le\exp(C(d_1+\log(1/b)))$. For the final inequality
use $\alpha\ge e^{-d}$ on a live plateau and the definition of $d_1$ on a
static plateau.

For the transfer from one node to another, the phase of any signed sum
$v\in[D]$ is approximated by the corresponding product of rounded phases
with uniform error at most $c_5b^6$. Scalar moments consequently change by
$O((c_5+c_6)b^6)$ between a prospective child measure and the stored post
measure if no refresh occurs. For the main functional, if $\mu$ is a joint
mass measure of input phase and side value, write $M$ for its side marginal
and $f$ for its phase moment. Then
\begin{equation}\label{eq:07-bounded-test}
 \int\frac{|f|^2}{M}
 =\sup_{|\gamma(z)|\le1}
       \int\bigl(2\operatorname{Re}(\gamma(z)\chi_v)-|\gamma(z)|^2\bigr)
                         \,d\mu.
\end{equation}
The tests have absolute value at most three. Phase rounding contributes at
most twice the phase error times the total mass, which is at most four;
total variation contributes at most three times the distance. Hence the
same $O((c_5+c_6)b^6)$ transfer applies to the main functional. A prospective
child high at $0.9b$ therefore makes the stored post law high with ample
slack above $b/2$. The scalar triangle inequality and the convexity of
$(M,f)\mapsto |f|^2/M$ imply that at least one prospective child of that
stored post law is high at $b/2$. Its signed sum belongs to the stored
cover. If a fresh snapshot is born at the current node, use its prospective
cover directly. This proves the transfer assertion in every case.
\end{proof}

\begin{proof}[Proof of the birth count]
Each rounded coordinate is observed only once, when appended, and costs at
most $C\log(CD_{\max}/b)$ in conditional entropy. The successive coordinate
spaces increase by retaining all previous coordinates. Before cap breach
there are at most $D_{\max}$ new coordinates on a path. Thus
\Cref{lemma:07-local-budget} gives
\begin{equation}\label{eq:07-rounded-fisher}
 \mathbb E_{\mathbf U_e}\sum_p
       \int\frac{|D_{\rm jump}|^2}{M}
 \le C d_1+C D_{\max}\log(CD_{\max}/b).
\end{equation}
Here $M$ has the full coordinate list available at that jump. A Fisher term
computed after forgetting some coordinates is smaller, by conditional
Jensen. This permits comparisons in old coordinate spaces while charging
all of them to \Cref{eq:07-rounded-fisher}.

\Needspace{12\baselineskip}
Start a comparison interval at each snapshot's post state in its fixed
coordinates. If its next replacement is caused by an append, or if the cap
is breached, end the old comparison before that node's jump, in the old
coordinates. Intervening removals at that node can be included. If the next
replacement is caused by distance with no append, include its node's jump
in the \emph{old} comparison and end at the selected child pre measure.
The new snapshot nevertheless starts at that trigger node's post measure,
before this same bit. A terminal interval ends at the plateau stop.
\Cref{fig:07-snapshot-timing} records this order. Inclusion of a trigger
bit is predictable, because the trigger is a test of both prospective
children before either is selected.

\begin{figure}[tbp]
 \centering
 \begin{tikzpicture}[
  x=1cm,y=1cm,>=Latex,
  state/.style={draw=black!65,rounded corners=2pt,fill=black!3,
    text width=3cm,minimum height=1.85cm,align=center,font=\small},
  every node/.style={font=\small},
  comparison/.style={very thick,>=Latex}
]
  \node[draw=black!55,rounded corners=2pt,fill=black!3,
    text width=14.7cm,align=center,inner sep=6pt] at (6,5.1)
    {\textbf{Before $r_{j+1}$: test both prospective children against the old snapshot $M_*$.}\\[3pt]
     $M_p^\pm=M_p\pm D_p,\qquad
       \max_{\pm}\|M_p^\pm-M_*\|_{\rm var}>c_6b^6$.\\[3pt]
     The selected child is far with probability at least $1/2$.};

  \node[align=center,text width=3.6cm] at (0,3.25)
    {\textbf{Replace now:}\\store the post measure $M_p$};
  \node[state] (post) at (0,1.6)
    {Post state $p$ (depth $j$)\\[3pt]
     $M_p$ after removals\\[2pt]
     coefficient $c_p$ fixed};
  \node[state] (prechild) at (4,1.6)
    {Selected child pre\\[3pt]
     $M_p+r_{j+1}D_p$\\[2pt]
     before any removals};
  \node[state] (postchild) at (8,1.6)
    {Next prefix (depth $j+1$)\\[3pt]
     $\epsilon_j$ known\\[2pt]
     process child if not stopped};
  \node[state] (output) at (12,1.6)
    {Delayed output\\[3pt]
     $c_p\epsilon_j r_{j+2}$\\[2pt]
     one bit later};

  \draw[->,thick] (post.east) -- node[above] {$r_{j+1}$} (prechild.west);
  \draw[->,thick] (prechild.east) -- (postchild.west);
  \draw[->,thick] (postchild.east) -- node[above] {$r_{j+2}$} (output.west);

  \draw[dashed,black!55] (prechild.south) -- (4,0);
  \draw[comparison,->] (-1.3,0) -- (4,0);
  \node[align=center,text width=6.3cm,fill=white,inner sep=2pt] at (1.35,-0.65)
    {Old comparison includes $r_{j+1}$.\\Stop at the selected child pre measure.};

  \draw[comparison,->] (0,-1.5) -- (13.3,-1.5);
  \node[align=center,text width=13.2cm] at (6.6,-2.05)
    {New comparison starts at $M_p$ before the same bit and continues forward.\\
     Only the trigger bit overlaps; an append ends the old comparison before that bit.};
\end{tikzpicture}

 \caption{A distance trigger without an append. At the post state $p$, both
 child pre measures are compared with the old snapshot $M_*$. The snapshot
 is replaced by $M_p$ before $r_{j+1}$ is read. The old comparison includes
 that bit and ends at the actual child pre measure; the new comparison
 starts at $M_p$, so this bit is used twice at most. Any child removals occur
 afterward. The coefficient of the jump from $p$ is emitted using the
 delayed sign $r_{j+2}$, even if no child jump is made. Here $c_p$ denotes
 that fixed normalized transform coefficient. An append instead ends the old comparison before
 $r_{j+1}$, in its old coordinates.}
 \label{fig:07-snapshot-timing}
\end{figure}

On one comparison interval let $k$ be each removed positive joint mass
measure, expressed in its fixed coordinate space. Retain $k$ unchanged
after its removal and let $K$ be the sum retained so far. Then
\[
 N=M+K
\]
is a positive measure martingale under the uniform output bits. At a kill,
$M$ loses precisely the measure added to $K$. At a spatial bit its two
values are $N\pm D_{\rm jump}$, with
$|D_{\rm jump}|\le M\le N$. In particular the martingale starts from the
snapshot measure $M_*$, and all later measures are absolutely continuous
with respect to $M_*$: a zero coordinate set remains zero under positive
child averaging and removals.

Use generalized relative entropy for finite measures,
\[
 \mathcal E(N\mid M_*)=
       \int(n\log n-n+1)\,dM_*,\qquad n=dN/dM_*.
\]
At a spatial bit, writing $t=D_{\rm jump}/N$, its average increment is
\[
 \int \frac N2\bigl((1+t)\log(1+t)+(1-t)\log(1-t)\bigr)
 \le C\int\frac{|D_{\rm jump}|^2}{N}
 \le C\int\frac{|D_{\rm jump}|^2}{M}.
\]
The scalar ratio in this inequality is bounded on $-1\le t\le1$, with
its limit at zero understood. Kills leave $N$ unchanged. Bounded stopping
in the finite-depth tree therefore bounds the expected terminal entropy
by the expected Fisher sum on the comparison interval.

At a distance trigger, one of the two prospective children differs from
$M_*$ by more than $c_6b^6$, so the selected child has this property with
conditional probability at least one half. On that event either
$\|K\|\ge c_6b^6/2$, or
$\|N-M_*\|_{\rm var}>c_6b^6/2$. On the complementary event $\|K\|<c_6b^6/2$, the latter alternative
holds, $N$ has mass at most $4+c_6b^6/2$, and $M_*$ has mass at most two. The scalar inequality
\[
 \frac{(x-1)^2}{1+x}\le C(x\log x-x+1),\qquad x\ge0,
\]
followed by Cauchy--Schwarz gives
\[
 \|N-M_*\|_{\rm var}^2
 \le C\bigl(N(\Omega)+M_*(\Omega)\bigr)\mathcal E(N\mid M_*),
\]
where $\Omega$ is the fixed observation space of the comparison interval.
Thus the latter event forces
$\mathcal E(N\mid M_*)\ge c b^{12}$. Conditional on the comparison start,
the probability that its next replacement is a distance trigger is at most
\begin{equation}\label{eq:07-trigger-charge}
 C\mathbb E\left[
       b^{-6}\sum_{\rm interval}\xi_p
       +b^{-12}\sum_{\rm interval}
                      \int\frac{|D_{\rm jump}|^2}{M}
             \,\middle|\,\text{comparison start}\right].
\end{equation}
Indeed the far-child event has at least half the trigger probability, and
the preceding deterministic alternatives charge it to removals or to
terminal entropy, whose expectation was just bounded.

Sum \Cref{eq:07-trigger-charge} with the probabilities of the starts.
Every jump is in at most two comparison intervals: at a distance trigger
it is the final bit of the old interval and the first bit of the new one.
Appends stop old comparisons before their jump. Removals likewise have
bounded overlap. By \Cref{eq:07-kill-bounds,eq:07-rounded-fisher}, distance
triggers contribute the last two terms in
\Cref{eq:07-snapshot-count}. Initialization contributes at most one, and
append replacements at most $D_{\max}$. This proves the count. No coordinate
is re-observed merely because a snapshot is replaced, and the conditional
argument did not condition on a realized side or rounded value.
\end{proof}

\subsection{Activation and the one block of missed jumps}

\begin{lemma}[Activation and localization before activation]
\label{lemma:07-activation}
On a capped scheme, declare a high $u$ activated at a jump $p$ if there is
$v\in[D]$ such that $|u-v||p|\le\sigma$ in the main case, or
$|u-v||p|\le c_8b^4$ in the scalar case. A representative from the current
snapshot then gives the fixed decompositions and bounds in
\Cref{prop:capture}. All non-bad high jumps before first activation lie in
one block of at most $H_b$ depths as stated there.
\end{lemma}

\begin{proof}
Freeze the spatial phase $\chi_{u-v}$ at any point of $p$. The normalized
scalar error on a prospective child is at most $C|u-v||p|$, by its mass
bound. In the main case conditional Cauchy--Schwarz gives
\[
 \left|\sqrt{G_{p'}^u}-\sqrt{G_{p'}^v}\right|
 \le\left(\frac1\alpha\mathbb E_{p'}
       [a|\chi_{u-v}(Y)-\chi_{u-v}(y_p)|^2]\right)^{1/2}
 \le C|u-v||p|.
\]
Thus $v$ is high at $0.9b$, since $\sigma\ll b_0$ and the scalar
prefactor is small. By \Cref{lemma:07-snapshots}, it is covered by a stored
$s$. If that snapshot was born at $q\supseteq p$, its main cover gives
$v-s=n+j$, with $\Delta_A(n)\le r_{\rm sep}$ and $|j||q|\le\sigma$.
Set $a'=u-v+j$. Then $|a'||p|\le2\sigma$. The channel $A$ is fixed
throughout the plateau, so $n$ retains its metric bound. The scalar cover
instead gives $|v-s||q|\le c_7b^4$, and hence
$|(u-v)+(v-s)||p|\le(c_7+c_8)b^4$. Fix this $s,n,a'$ at activation;
all their spatial error bounds only improve on descendants.

Now let $p_0$ be the first non-bad high jump before any activation, if one
exists. By \Cref{eq:source-14} there is a signed combination $v\in[D]$ at
that node with
\[
 |u-v||p_0|\le P(D_{\max},d_1,1/b)
\]
for a fixed polynomial $P$. The same signed combination remains available
as the list grows. After $r$ further depths the left side with the new
interval is at most $2^{-r}P(D_{\max},d_1,1/b)$. For
\[
 r\ge C\log(CD_{\max}d_1/b)+C\log(1/\sigma)
\]
this is at most the relevant activation tolerance, $\sigma$ or $c_8b^4$.
Any subsequent high jump then activates, whether or not it is bad. Hence
every non-bad high jump strictly before first activation lies in the one
depth block beginning at $p_0$ of the claimed length. If no such $p_0$
exists there is no missed non-bad high jump. This proves the assertion and
completes \Cref{prop:capture}.
\end{proof}

The cap has been deterministic throughout the probability estimates. In
\Cref{sec:paths} the estimates will first be made simultaneous over all
dyadic caps, after which a cap comparable to $1+|D|_{\rm final}$ may be
selected on each path. In particular \Cref{eq:07-append-budget} and Jensen
give
\[
 \mathbb E_{\mathbf U_e}\log(C(1+|D|_{\rm final}))
       \le C\log(Cd_1/b),
\]
so this later choice incurs a logarithmic expected cost and excludes no
large-list paths.
